\section{Propriet\`a dello Schema dell'Ontologia}
\label{sec:tbox_proprieta}

\subsection{Object Property}

La Tabella~\ref{tab:objprop} elenca tutte le 31 object property definite nell'ontologia.
Per semplicit\`a, si omettono le propriet\`a inverse quando speculari; la colonna
\textbf{Assiomi} indica le caratteristiche OWL 2 aggiuntive.

\begin{longtable}{lllll}
\caption{Object Property}\label{tab:objprop}\\
\toprule
\textbf{Propriet\`a}      & \textbf{Domain} & \textbf{Range}   & \textbf{Assiomi}                 & \textbf{Label} \\
\midrule
\endfirsthead
\multicolumn{5}{c}{\tablename\ \thetable{} -- continua} \\
\toprule
\textbf{Propriet\`a}      & \textbf{Domain} & \textbf{Range}   & \textbf{Assiomi}                 & \textbf{Label} \\
\midrule
\endhead
\bottomrule
\endlastfoot

developedBy        & VideoGame & Developer   & Asym, Irr, inv developerOf            & developed by \\
developerOf        & Developer & VideoGame   & Asym, Irr                             & developer of \\
publishedBy        & VideoGame & Publisher   & Asym, Irr, inv publisherOf            & published by \\
publisherOf        & Publisher & VideoGame   & Asym, Irr                             & publisher of \\
hasGenre           & VideoGame & Genre       & Asym, Irr, inv genreOf                & has genre \\
genreOf            & Genre     & VideoGame   & Asym, Irr                             & genre of \\
availableOn        & VideoGame & Platform    & Irr, inv platformFor                  & available on \\
platformFor        & Platform  & VideoGame   & Asym, Irr                             & platform for \\
hasCharacter       & VideoGame & Character   & Asym, Irr, inv appearsIn              & has character \\
appearsIn          & Character & VideoGame   & Asym, Irr                             & appears in \\
belongsTo          & VideoGame & Franchise   & Asym, Irr, inv includes               & belongs to \\
includes           & Franchise & VideoGame   & Asym, Irr                             & includes \\
wonAward           & VideoGame & Award       & Irr, inv awardWonBy                   & won award \\
awardWonBy         & Award     & VideoGame   & Irr                                   & award won by \\
sequelOf           & VideoGame & VideoGame   & Asym, Irr, Trans, inv hasSequel       & sequel of \\
hasSequel          & VideoGame & VideoGame   & Asym, Irr                             & has sequel \\
hasGameMode        & VideoGame & GameMode    & Irr, inv modeInGame                   & has game mode \\
modeInGame         & GameMode  & VideoGame   & Irr                                   & mode in game \\
madeWith           & VideoGame & GameEngine  & Asym, Irr, inv engineUsedIn           & made with \\
engineUsedIn       & GameEngine& VideoGame   & Asym, Irr                             & engine used in \\
partOfFranchise    & Franchise & Franchise   & Asym, Irr, Trans                     & part of franchise \\
subsidiaryOf       & ---       & ---         & Asym, Irr, Trans                     & subsidiary of \\
sharedFranchiseWith& VideoGame & VideoGame   & Sym, chain: belongsTo $\circ$ includes & shared franchise with \\
sharesDeveloperWith& VideoGame & VideoGame   & Sym, chain: developedBy $\circ$ developerOf & shares developer with \\
sharesPublisherWith& VideoGame & VideoGame   & Sym, chain: publishedBy $\circ$ publisherOf & shares publisher with \\
reviewer           & Review    & ---         & ---                                   & reviewer \\
reviewedGame       & Review    & VideoGame   & Funct, inv hasReview                  & reviewed game \\
hasReview          & VideoGame & Review      & ---                                   & has review \\
reviewOf           & Review    & VideoGame   & ---                                   & review of \\
eventFor           & GameEvent & VideoGame   & inv hasEvent                          & event for \\
hasEvent           & VideoGame & GameEvent   & ---                                   & has event \\
\bottomrule
\end{longtable}

\textit{Legenda}: Asym = AsymmetricProperty, Irr = IrreflexiveProperty,
Sym = SymmetricProperty, Trans = TransitiveProperty, Funct = FunctionalProperty,
inv = inverseOf, chain = propertyChainAxiom.

\subsection{Datatype Property}

La Tabella~\ref{tab:dataprop} elenca tutte le 26 datatype property.

\begin{longtable}{lllll}
\caption{Datatype Property}\label{tab:dataprop}\\
\toprule
\textbf{Propriet\`a}      & \textbf{Domain} & \textbf{Range}   & \textbf{subPropertyOf} & \textbf{Label} \\
\midrule
\endfirsthead
\multicolumn{5}{c}{\tablename\ \thetable{} -- continua} \\
\toprule
\textbf{Propriet\`a}      & \textbf{Domain} & \textbf{Range}   & \textbf{subPropertyOf} & \textbf{Label} \\
\midrule
\endhead
\bottomrule
\endlastfoot

name              & ---         & xsd:string  & ---     & name \\
gameName          & VideoGame   & xsd:string  & name    & game name \\
gameDescription   & VideoGame   & xsd:string  & ---     & game description \\
releaseDate       & VideoGame   & xsd:date    & ---     & release date \\
metacriticScore   & VideoGame   & xsd:integer & ---     & Metacritic score \\
developerName     & Developer   & xsd:string  & name    & developer name \\
publisherName     & Publisher   & xsd:string  & name    & publisher name \\
genreName         & Genre       & xsd:string  & name    & genre name \\
platformName      & Platform    & xsd:string  & name    & platform name \\
characterName     & Character   & xsd:string  & name    & character name \\
franchiseName     & Franchise   & xsd:string  & name    & franchise name \\
awardName         & Award       & xsd:string  & name    & award name \\
awardYear         & Award       & xsd:integer & ---     & award year \\
description       & ---         & xsd:string  & ---     & description \\
officialWebsite   & VideoGame   & xsd:anyURI  & ---     & official website \\
countryOfOrigin   & VideoGame   & xsd:string  & ---     & country of origin \\
engineName        & GameEngine  & xsd:string  & name    & engine name \\
hasSteamAppId     & VideoGame   & xsd:string  & ---     & Steam App ID \\
hasWikidataId     & ---         & xsd:string  & ---     & Wikidata ID \\
hasIGDBId         & VideoGame   & xsd:string  & ---     & IGDB ID \\
reviewScore       & Review      & xsd:integer & ---     & review score \\
reviewDate        & Review      & xsd:date    & ---     & review date \\
reviewerName      & Review      & xsd:string  & ---     & reviewer name \\
eventDate         & GameEvent   & xsd:date    & ---     & event date \\
budgetTier        & VideoGame   & xsd:string  & ---     & budget tier \\
playerCount       & VideoGame   & xsd:integer & ---     & player count \\
\bottomrule
\end{longtable}

\subsection{Property Chain Axiom}

L'ontologia definisce tre \texttt{owl:propertyChainAxiom}, che permettono di
inferire relazioni transitive composte. In notazione DL (Description Logic):

\begin{enumerate}[nosep]
  \item \texttt{sharedFranchiseWith} $\equiv$ \texttt{belongsTo} $\circ$ \texttt{includes}

    Due giochi condividono un franchise se il primo appartiene a un franchise
    e il franchise include il secondo.
  \item \texttt{sharesDeveloperWith} $\equiv$ \texttt{developedBy} $\circ$ \texttt{developerOf}

    Due giochi condividono lo sviluppatore se il primo \`e sviluppato da uno
    studio che ha sviluppato anche il secondo.
  \item \texttt{sharesPublisherWith} $\equiv$ \texttt{publishedBy} $\circ$ \texttt{publisherOf}

    Analogo per l'editore.
\end{enumerate}

Queste catene sono materializzate dal reasoner OWL 2 RL (\texttt{owlrl})
e producono relazioni simmetriche tra giochi dello stesso franchise, studio o editore.
Per evitare esplosioni combinatorie su franchise/studio/editore con migliaia di
giochi, il materialiser mirato in \texttt{enrich\_owl.py} limita il numero di
coppie per gruppo tramite il flag \texttt{--max-group} (default 50).

\subsection{Propriet\`a Transitive}

Due propriet\`a sono dichiarate transitive:

\begin{itemize}[nosep]
  \item \texttt{sequelOf} --- se $A$ \`e sequel di $B$ e $B$ \`e sequel di $C$,
    allora $A$ \`e sequel di $C$. Permette di tracciare catene di sequel
    (es. Dark Souls III $\rightarrow$ Dark Souls II $\rightarrow$ Dark Souls);
  \item \texttt{partOfFranchise} --- per gerarchie di franchise con sotto-franchise.
    Es. \textless{}Mario Kart\textgreater{} $\rightarrow$ \textless{}Mario\textgreater{}.
\end{itemize}

La transitivit\`a di \texttt{sequelOf} \`e particolarmente utile per query del
tipo \textless{}tutti i sequel di Dark Souls\textgreater{} senza dover navigare
manualmente la catena.

\subsection{Coppie Inverse Sistematiche}

L'ontologia definisce 13 coppie inverse, consentendo la navigazione bidirezionale
del grafo. Per ogni propriet\`a ``forward'' esiste la corrispondente inversa:
\begin{itemize}[nosep]
  \item \texttt{developedBy} $\leftrightarrow$ \texttt{developerOf}
  \item \texttt{publishedBy} $\leftrightarrow$ \texttt{publisherOf}
  \item \texttt{hasGenre} $\leftrightarrow$ \texttt{genreOf}
  \item \texttt{availableOn} $\leftrightarrow$ \texttt{platformFor}
  \item \texttt{hasCharacter} $\leftrightarrow$ \texttt{appearsIn}
  \item \texttt{belongsTo} $\leftrightarrow$ \texttt{includes}
  \item \texttt{wonAward} $\leftrightarrow$ \texttt{awardWonBy}
  \item \texttt{sequelOf} $\leftrightarrow$ \texttt{hasSequel}
  \item \texttt{hasGameMode} $\leftrightarrow$ \texttt{modeInGame}
  \item \texttt{madeWith} $\leftrightarrow$ \texttt{engineUsedIn}
  \item \texttt{reviewedGame} $\leftrightarrow$ \texttt{hasReview}
  \item \texttt{eventFor} $\leftrightarrow$ \texttt{hasEvent}
\end{itemize}

La tredicesima coppia (\texttt{hasGenre} $\leftrightarrow$ \texttt{genreOf}) era
gi\`a implicita nello schema originale. La dichiarazione esplicita di \texttt{owl:inverseOf}
in una sola direzione \`e sufficiente: \texttt{owlrl} inferisce automaticamente
la direzione opposta durante la materializzazione. Il validatore dello schema segnala
12 warning \texttt{INVERSE\_NOT\_BIDIRECTIONAL} proprio per questa asimmetria
dichiarativa (cfr.~Sezione~\ref{sec:validazione}).
